\documentclass[../../main.tex]{subfiles}
\begin{document}
\section{de Rham 定理}\label{section:ch18-the-de-rham-theorem}

本节我们叙述并证明 de Rham 定理。在进入定理之前，我们还需要一个代数引理，其证明是另一次像蛇引理那样的图表追踪。

\begin{lemma}[五引理]\label{lemma:ch18-five-lemma}
考虑下面由模与线性映射构成的交换图，其中两行正合：
\begin{align*}
\begin{array}{ccccccccc}
A_1&\to&A_2&\to&A_3&\to&A_4&\to&A_5\\
\downarrow f_1&&\downarrow f_2&&\downarrow f_3&&\downarrow f_4&&\downarrow f_5\\
B_1&\to&B_2&\to&B_3&\to&B_4&\to&B_5
\end{array}
\end{align*}
若 $f_1,f_2,f_4,f_5$ 都是同构，则 $f_3$ 也是同构。
\end{lemma}

\begin{exercise}\label{exercise:ch18-11}
证明（或查阅）五引理。
\end{exercise}

设 $M$ 是一个光滑流形，$\omega$ 是 $M$ 上的闭 $p$-形式，$\sigma$ 是 $M$ 中的一个光滑 $p$-单形。我们定义 $\omega$ 在 $\sigma$ 上的积分为
\begin{align*}
\int_\sigma\omega=\int_{\Delta^p}\sigma^*\omega.
\end{align*}
如果 $c=\sum_i c_i\sigma_i$ 是一个光滑 $p$-链，定义 $\int_c\omega=\sum_i c_i\int_{\sigma_i}\omega$。

\begin{theorem}[链上的 Stokes 定理]\label{theorem:ch18-stokes-for-chains}
若 $c$ 是光滑流形 $M$ 中的一个光滑 $p$-链，$\omega$ 是 $M$ 上的光滑 $(p-1)$-形式，则
\begin{align}
\int_{\partial c}\omega=\int_c\mathrm{d}\omega.
\end{align}
\end{theorem}

\begin{proof}
只需对光滑单形 $\sigma$ 证明。由于 $\Delta^p$ 是带角流形，Stokes 定理给出 $\int_\sigma\mathrm{d}\omega=\int_{\partial\Delta^p}\sigma^*\omega$。各面映射 $F_{i,p}$ 给出 $\partial\Delta^p$ 的参数化，带符号 $(-1)^i$，而这恰是奇异边界算子的定义。
\end{proof}

利用这个定理，我们定义一个自然线性映射 $I:H^p_{\mathrm{dR}}(M)\to H^p(M;\mathbb{R})$，称为\emph{de Rham 同态}。对任意 $[\omega]\in H^p_{\mathrm{dR}}(M)$ 与 $[c]\in H_p(M)\cong H_p^\infty(M)$，定义
\begin{align}
I[\omega][c]=\int_{\widetilde c}\omega,
\end{align}
其中 $\widetilde c$ 是代表同调类 $[c]$ 的任意光滑 $p$-闭链。这是良定义的：两个代表同一类的闭链相差一个边缘，而闭形式在边缘上的积分为零；若 $\omega=\mathrm{d}\eta$ 恰当，则它在闭链上的积分也为零。

\begin{proposition}[de Rham 同态的自然性]\label{proposition:ch18-naturality-of-derham}
de Rham 同态 $I$ 与光滑映射诱导的上同调映射交换，也与 Mayer--Vietoris 序列中的连接同态交换。
\end{proposition}

\begin{figure}[ht]
\centering
\includegraphics[width=0.6\textwidth]{fig-18-7.png}
\caption{$I$ 关于连接同态的自然性}\label{figure:ch18-naturality-connecting}
\end{figure}

\begin{theorem}[de Rham]\label{theorem:ch18-derham}
对每个光滑流形 $M$ 与非负整数 $p$，de Rham 同态 $I:H^p_{\mathrm{dR}}(M)\to H^p(M;\mathbb{R})$ 是一个同构。
\end{theorem}

\begin{proof}
称一个光滑流形 $M$ 为一个 de Rham 流形，如果对每个 $p$，$I$ 都是同构。
\textbf{第一步：}可数个 de Rham 流形的不相交并仍是 de Rham 流形。
\textbf{第二步：}$\mathbb{R}^n$ 的每个凸开子集都是 de Rham 流形（由 Poincaré 引理与单点空间的上同调）。
\textbf{第三步：}若 $M$ 有一个有限 de Rham 覆盖，则 $M$ 是 de Rham 流形。这是证明的核心：对覆盖中集合的个数作归纳；对两个集合的覆盖，把 de Rham 与奇异上同调的 Mayer--Vietoris 序列拼成交换图，用五引理得出中间的 $I$ 也是同构。
\textbf{第四步：}若 $M$ 有一个 de Rham 基，则 $M$ 是 de Rham 流形。用一个穷竭函数把 $M$ 切成一层层，再用奇偶两层的并把问题归结到有限覆盖情形。
\textbf{第五步：}$\mathbb{R}^n$ 的每个开子集都是 de Rham 流形（有由凸球组成的基）。
\textbf{第六步：}每个光滑流形都是 de Rham 流形（有由光滑坐标域组成的基）。
\end{proof}

\begin{figure}[ht]
\centering
\includegraphics[width=0.8\textwidth]{fig-18-8.png}
\caption{de Rham 定理证明的第四步}\label{figure:ch18-derham-proof-step4}
\end{figure}

这个结果表达了光滑流形的拓扑性质与分析性质之间的深刻联系，在微分几何中处于中心地位。de Rham 定理有许多种证明；这里给出的优雅证明属于 Glen E.\ Bredon。
\end{document}
